%! TEX root = master.tex
\chapter{2}{The Galilean Group}

To turn the principle of relativity into mathematics, we need a precise
model of space-time and of the transformations that preserve its
structure \cite[ch.~1]{arnold}.

\section{Affine space}\label{sec:affine}

\begin{definition}[Affine space]\label{def:affine}
  \emph{Affine $n$-space} $A^n$ is a set of points on which the vector space
  $\R^n$ acts by parallel displacement. To a point $a \in A^n$ and a vector
  $\boldsymbol\xi \in \R^n$ corresponds the point $a + \boldsymbol\xi$, and
  for any two points $a, b$ there is exactly one vector
  $b - a \in \R^n$ carrying $a$ to $b$.
\end{definition}

\begin{remark}
  $A^n$ differs from $\R^n$ in having no distinguished origin. Points can be
  subtracted, giving a vector, but not added. Choosing an origin
  $o \in A^n$ identifies $A^n$ with $\R^n$ through $a \mapsto a - o$, but no
  choice is preferred. This is exactly the situation in physical space.
\end{remark}

\section{Galilean structure}\label{sec:galilean-structure}

\begin{definition}[Galilean space-time]\label{def:galilean-structure}
  A \emph{Galilean structure} consists of three things.
  \begin{enumerate}
    \item The \emph{universe}, a four-dimensional affine space $A^4$, whose
      points are called \emph{events}.
    \item \emph{Time}, a linear map $t : \R^4 \to \R$ from the displacements
      of the universe onto the real line. The \emph{time interval} from event
      $a$ to event $b$ is $t(b - a)$. Events with $t(b - a) = 0$ are
      \emph{simultaneous}.
    \item A \emph{distance} between simultaneous events,
      $\rho(a, b) = \sqrt{(a - b, a - b)}$, given by a scalar product on the
      three-dimensional kernel of $t$. With it, each set of simultaneous
      events is a Euclidean space $E^3$.
  \end{enumerate}
  The coordinate model is $\R \times \R^3$, with time $t$ and the standard
  distance on each slice $\set{t} \times \R^3$.
\end{definition}

\begin{definition}[Galilean group]\label{def:galilean-group}
  A \emph{Galilean transformation} is an affine map of $A^4$ to itself that
  preserves time intervals and distances between simultaneous events. These
  transformations form the \emph{Galilean group}.
\end{definition}

\begin{exercise}\label{ex:galilean-examples}
  Check that the following maps of $\R \times \R^3$ are Galilean:
  \begin{enumerate}[label=(\alph*)]
    \item a \emph{uniform motion} with velocity $\mathbf v$,
      $g_1(t, \mathbf x) = (t, \mathbf x + \mathbf v t)$;
    \item a \emph{translation},
      $g_2(t, \mathbf x) = (t + s, \mathbf x + \mathbf s)$;
    \item a \emph{rotation}, $g_3(t, \mathbf x) = (t, G\mathbf x)$, where
      $G : \R^3 \to \R^3$ is orthogonal.
  \end{enumerate}
\end{exercise}

\begin{solution}
  All three are affine. Take events $(t, \mathbf x)$ and $(t', \mathbf x')$.
  \begin{enumerate}[label=(\alph*)]
    \item The time interval is unchanged. If $t = t'$, the spatial difference
      $(\mathbf x' + \mathbf v t) - (\mathbf x + \mathbf v t) = \mathbf x' - \mathbf x$
      is unchanged too.
    \item Both times and both positions are shifted by the same amounts, so
      differences are unchanged.
    \item Times are unchanged, and
      $\abs{G\mathbf x' - G\mathbf x} = \abs{G(\mathbf x' - \mathbf x)} = \abs{\mathbf x' - \mathbf x}$
      because $G$ is orthogonal. \qedhere
  \end{enumerate}
\end{solution}

%! TEX root = ../master.tex
\begin{theorem}[Structure of the Galilean group]\label{thm:galilean-decomposition}
  Every Galilean transformation $g$ of $\R \times \R^3$ can be written
  uniquely as
  \[
    g = g_1 \circ g_2 \circ g_3,
  \]
  where $g_1$ is a uniform motion, $g_2$ a translation and $g_3$ a rotation
  (\cref{ex:galilean-examples}).
  \thmsource{\cite[ch.~1]{arnold}}
\end{theorem}

\begin{proof}
  A Galilean transformation is affine, so $g(t, \mathbf x) = L(t, \mathbf x) + (s, \mathbf s)$
  for a linear map $L$ and a fixed $(s, \mathbf s) \in \R \times \R^3$.

  \emph{Time.} $g$ preserves time intervals, so the time component of
  $L(\tau, \boldsymbol\xi)$ is $\tau$.

  \emph{Space.} Hence $L$ sends a simultaneous displacement
  $(0, \boldsymbol\xi)$ to a simultaneous displacement $(0, G\boldsymbol\xi)$
  for some linear $G : \R^3 \to \R^3$. Distances between simultaneous events
  are preserved, so $\abs{G\boldsymbol\xi} = \abs{\boldsymbol\xi}$, and $G$ is
  orthogonal.

  \emph{Velocity.} Write $L(1, \mathbf 0) = (1, \mathbf v)$. By linearity
  $L(\tau, \boldsymbol\xi) = (\tau, \mathbf v\tau + G\boldsymbol\xi)$, so
  \[
    g(t, \mathbf x) = (t + s,\; G\mathbf x + \mathbf v t + \mathbf s).
  \]
  \emph{Factorisation.} For $g_1(t, \mathbf x) = (t, \mathbf x + \mathbf v t)$,
  $g_2(t, \mathbf x) = (t + s, \mathbf x + \mathbf u)$ and
  $g_3(t, \mathbf x) = (t, G\mathbf x)$,
  \[
    g_1 g_2 g_3 (t, \mathbf x) = \bigl(t + s,\; G\mathbf x + \mathbf u + \mathbf v(t + s)\bigr).
  \]
  This equals $g$ exactly when $\mathbf u = \mathbf s - \mathbf v s$.

  \emph{Uniqueness.} The composite determines $s$ (its time shift), $G$
  (its action on $\mathbf x$), $\mathbf v$ (its coefficient of $t$) and then
  $\mathbf u$. So two factorisations of the same $g$ coincide.
\end{proof}


\begin{remark}
  Counting parameters: $\mathbf v$ gives $3$, $(s, \mathbf s)$ gives $4$,
  and the orthogonal $G$ gives $3$ (three angles, plus a choice of sign for
  $\det G$). So the Galilean group is $10$-dimensional.
\end{remark}

\section{Relativity in the equations}\label{sec:relativity-equations}

In this language, Galileo's principle says that Newton's equation
(\cref{thm:newton-equation}) is unchanged by Galilean transformations: if
$g$ is Galilean and $\mathbf x(t)$ is a motion, so is its image under $g$.
Each part of the group then forces a property of $\mathbf F$.

\begin{exercise}\label{ex:time-homogeneity}
  Use invariance under the time translations $(t, \mathbf x) \mapsto (t + s, \mathbf x)$
  to show that $\mathbf F$ does not depend on $t$.
\end{exercise}

\begin{proof}
  Let $\mathbf x(t)$ be a motion. Its image under the translation is the
  motion $\mathbf y(t) = \mathbf x(t - s)$, which by assumption satisfies the
  same equation:
  \[
    \ddot{\mathbf y}(t) = \mathbf F\bigl(\mathbf y(t), \dot{\mathbf y}(t), t\bigr).
  \]
  But $\ddot{\mathbf y}(t) = \ddot{\mathbf x}(t - s)
  = \mathbf F\bigl(\mathbf x(t - s), \dot{\mathbf x}(t - s), t - s\bigr)$.
  Comparing the two,
  \[
    \mathbf F(\mathbf x_0, \mathbf v_0, t) = \mathbf F(\mathbf x_0, \mathbf v_0, t - s),
  \]
  where $\mathbf x_0 = \mathbf x(t - s)$ and $\mathbf v_0 = \dot{\mathbf x}(t - s)$.
  By determinacy there is a motion through any initial state
  $(\mathbf x_0, \mathbf v_0)$ at any time. So the identity holds for all
  arguments and all $s$, and $\mathbf F = \mathbf F(\mathbf x, \dot{\mathbf x})$.
\end{proof}

\begin{remark}
  In the same way, invariance under spatial translations forces $\mathbf F$ to
  depend only on the relative positions $\mathbf x_i - \mathbf x_j$. Uniform
  motions force it to depend only on the relative velocities. Rotations make
  it equivariant: rotating all positions and velocities rotates the
  accelerations.
\end{remark}
